\section{Introduction}
\label{sec:intro}

Sorting waste at the point of disposal, for example in a camera-assisted bin or collection station that suggests a category for the item presented to it, needs a classifier that runs locally on inexpensive hardware, with a small memory footprint and low latency.
The data for such a task are small and imbalanced. TrashNet~\citep{yang2016trashnet}, the dataset we use, has six material classes whose training counts range from 96 to 416 images in our split.
Compact convolutional networks such as MobileNets~\citep{sandler2018mobilenetv2,howard2019searching} and 8-bit integer inference~\citep{jacob2018quantization,krishnamoorthi2018quantizing} address the footprint, at some cost in accuracy relative to large networks.

Knowledge distillation (KD)~\citep{hinton2015distilling} recovers part of that accuracy. The compact student is trained to match a larger teacher's output distribution in addition to the hard labels.
Relational and feature-based variants transfer the geometry of the teacher's representation instead of, or in addition to, its outputs~\citep{romero2015fitnets,park2019relational,tian2020contrastive}.
Only the student is deployed, so distillation is compatible with quantization, and the two are complementary compression mechanisms~\citep{polino2018model,kim2019qkd}.

A recent line of work asks whether quantum computation can assist distillation while the deployed model stays classical.
Quantum relational knowledge distillation (QRKD)~\citep{liu2025qrkd} extends relational KD. It maps features into a Hilbert space, treats them as quantum states, and uses quantum-kernel values to describe relations between samples. It reports gains over classical relational KD on CNNs for MNIST and CIFAR-10 and on GPT-2 language tasks, with quantum computation needed only in training.
Other work distills classical teachers into quantum students~\citep{li2024hybrid,hasan2023bridging} or uses a quantum model as a soft-label teacher for a classical student~\citep{naqvi2026structured}.
Our question is narrower: does the quantum component contribute anything that a classical kernel in the same place would not?

For kernels, that question has no obvious answer.
Quantum kernels can concentrate, so that kernel values stop distinguishing inputs~\citep{thanasilp2024exponential}. Numerical studies of fidelity kernels on classical data find no advantage~\citep{slattery2023numerical}. The inductive bias of a quantum kernel determines what it can learn~\citep{kubler2021inductive}, and once the bandwidth is tuned it can resemble a classical kernel~\citep{florezablan2025similarity}.
Large-scale benchmarks of fidelity and projected quantum kernels examine which design choices make them work~\citep{schnabel2025scrutiny}.
The projected quantum kernel (PQK) of Huang et al.~\citep{huang2021power} compares inputs through single-qubit reduced states of an embedding circuit. It is a clean candidate to test: at the small scale used here it can be simulated classically, and a classical kernel acting on the same inputs can replace it one-for-one.

Two practical points shape how such a comparison has to be run.
First, effects on small datasets are comparable in size to seed-to-seed noise. Conclusions therefore need several seeds, paired comparisons, and an explicit statement of what was selected on which split.
Second, the end point is the deployed INT8 model, and INT8 results depend on the kernels that execute it.

We therefore compare quantum-assisted distillation with classical alternatives under one pipeline.
TrashNet (six classes, imbalanced, macro-F1) is the primary dataset. RPS-25, a 25\% training subset of Rock--Paper--Scissors~\citep{moroney2019rps} (three classes), is a low-data confirmation.
A frozen ResNet-50 teacher~\citep{he2016deep} distills into three students of 0.03\,M to 0.93\,M parameters (LeNet-5, MobileNetV2-0.35, MobileNetV3-Small). Five methods are compared: hard-label training (\Scratch), logit KD (\Kd), relational KD (\Rkd), classical RBF-kernel alignment (\Rbf, the matched twin) and PQK (\Pqk).
All methods on a pair share data, optimizer and epoch budget. Each cell has \NSeeds{} seeds, and every exported model is a plain INT8 network with no quantum operation.
Two further arms isolate where a PQK effect could come from: the same circuit with its CNOT gates removed, and a fitted instead of random teacher projection.

Our findings are as follows.
\begin{itemize}
  \item The pipeline can detect a distillation gain: \Kd beats \Scratch on TrashNet for both MobileNets (by \KdGainTnA{} and \KdGainTnB{} macro-F1).
  \item \Pqk is not distinguishable from \Kd. It is ahead on the mean in \KdAheadFp{} of \NPairs{} pairs (\KdAheadFpEx{} if one seed trained with a different KD temperature is excluded) and behind on both TrashNet MobileNets (by \KdBehindTnA{} and \KdBehindTnB{}). The smallest uncorrected paired $p$-value is \MinPKd{}.
  \item At its pre-specified kernel scale, \Pqk beats its RBF twin in \GtwoFixedWins{} of \NPairs{} pairs (smallest uncorrected $p=\MinPRbfFixed$). With a median-heuristic scale, adopted after these results were seen, it does so in \GtwoAdaptWins{} of \GtwoAdaptPairs{}. The one nominally significant difference ($p=\MinPRbfAdapt$, RPS-25/MobileNetV3-Small) favours \MinPRbfAdaptFavours.
  \item QRKD reports consistent gains over classical relational KD. Our setting differs in data, students and loss. There, \Pqk is ahead of \Rkd in \RkdAhead{} of \RkdPairs{} pairs, by at most \RkdMaxLead{} macro-F1, and behind in the rest.
  \item On macro-F1, \Pqk has lower seed-to-seed variance than \Rbf in \VarBelowRbf{} of \NPairs{} pairs and than \Scratch in \VarBelowScratch{} (\VarBelowScratchEx{} if one recipe-mismatched run is excluded). On top-1 the counts are \VarTopBelowRbf{} and \VarTopBelowScratch{}. This is not significant at three seeds, and \Rkd is about as stable.
  \item Deployment does not depend on the training method: all \NIntModels{} exported MobileNet students are fully integer, per-channel INT8 models, and INT8 conversion changes macro-F1 by $\IntDeltaMean$ on average. The size of each method's INT8 effect depends on the runtime kernels.
\end{itemize}

All comparisons are exploratory. Cells have three seeds, and test splits informed earlier design choices. We therefore select on validation, report paired differences and treat test scores as descriptive.

\paragraph{Organization.}
Section~\ref{sec:related} reviews related work. Section~\ref{sec:method} describes the objectives, the quantum kernel and its controls, and Section~\ref{sec:setup} the data, training and deployment protocol.
Section~\ref{sec:results} reports the results\extra{, Section~\ref{sec:discussion} discusses them and Section~\ref{sec:limitations} lists their limits}.
\extra{The appendices document the problems we corrected and the details of the screens, the export and the on-device measurements.}

\paragraph{Contributions.}
(i)~A matched-control comparison of quantum-assisted and classical distillation for lightweight waste classification, with one pipeline, a classical RBF-kernel twin for the quantum kernel, and paired per-seed reporting.
(ii)~Two ablations that probe where a quantum-kernel effect could originate: removing the circuit's entanglement, and replacing the random teacher projection by a fitted one. The fitted projection makes the target kernel nearly the label kernel on training batches and did not help.
(iii)~Deployment evidence for the resulting students: a fully integer INT8 export of all \NIntModels{} models, on-device latency, memory and temperature on a Raspberry Pi~4B, and a measure of how INT8 results depend on the kernels that run them.
(iv)~An account of the pitfalls we hit and corrected, including a collapse of the kernel-alignment loss when the teacher-side projection is trainable, and of how our own selection choices limit what the results can support.
